\documentclass{article}

\usepackage[utf8]{inputenc}
\usepackage[margin=1in]{geometry}
\usepackage[titletoc,title]{appendix}

\usepackage{amsmath,amsfonts,amssymb,mathtools}
\usepackage{graphicx,float}
\usepackage{booktabs}
\usepackage[ruled,vlined]{algorithm2e}
\usepackage{algorithmic}
\usepackage{biblatex}
\usepackage{cleveref}
\usepackage{indentfirst}
\setlength{\parindent}{2em}

\addbibresource{references.bib}

\title{PRML Spring 2026 Homework 4: Transformer Models}
\author{Yifan Ding \\
        \small Code: \url{https://github.com/douhenhuidyf/PRML-Spring-2026/tree/main/hw4} \\
        \small Email: 23373493@buaa.edu.cn}
\date{May 15, 2026}

\begin{document}
\maketitle

% \begin{abstract}

% \end{abstract}


\section{Introduction}
\subsection{Problem Statement}
% 1）阅读并复现经典的17年文章： Attention is all you need!
% https://arxiv.org/abs/1706.03762

% （2）并尝试并通过实验来理解Transformer每个模块的必要性，可以通过实验来研究下面任何情况之一即可：

% 2.1 位置编码： 如果用简单的绝对编码或者其他位置编码形式来对比原文中的方法，讨论位置编码的核心作用是什么？
% 2.2 Q, K, V 的必要性： 如果我们只用Q与K （K, V合并为一个矩阵）对算法的影响是什么？
% 2.3 如果模块中不采用ResNet的结构，对于算法的影响是什么？
% 2.4 引入位置编码后用CNN重构一个网络，它与自注意力机制的Transformer比较
% 2.5 经典Transformer的模块哪个方向可以做更多的改进？（不是直接引用其他已发布论文的改进方法） 提出独立想法并实践验证你的想法可行性。
Read and reproduce the classic 2017 paper: "Attention is All You Need!"\cite{vaswani2023attentionneed}.
Then, through experiments, try to understand the necessity of each module in the Transformer.
You can choose to study any of the following situations:

1. Positional Encoding: Compare the original method with simple absolute encoding or other forms of positional encoding, and discuss the core role of positional encoding.

2. Necessity of Q, K, V: If we only use Q and K (merging K and V into one matrix), what is the impact on the algorithm?

3. If the ResNet structure is not adopted in the module, what is the impact on the algorithm?

4. After introducing positional encoding, reconstruct a network using CNN and compare it with the self-attention mechanism of the Transformer.

5. Which direction can be improved in the classic Transformer module? (Not directly citing other published improvement methods) Propose independent ideas and practice to verify the feasibility of your ideas.

\subsection{Problem Analysis}
The Transformer model, introduced by Vaswani et al. in 2017, was originally designed for natural language processing tasks, especially machine translation.
Thus, we will focus on the performance of the Transformer in translation tasks in this homework.
As training a Transformer from scratch can be computationally expensive, we will use a smaller Chinese to English translation dataset and a smaller scale configuration for our experiments.

\subsection{Structure of this paper}

\section{Methodology}
\label{sec:methodology}

\subsection{Metrics}
\label{sec:metrics}
To evaluate the performance of the baseline Transformer model and its variants, we use two standard machine translation metrics: BLEU and chrF. Both are computed at the corpus level.
\begin{itemize}
        \item \textbf{BLEU} (Bilingual Evaluation Understudy) measures $n$-gram precision with a brevity penalty to discourage overly short translations. It is defined as:
        \[
        \text{BLEU} = \text{BP} \cdot \exp\left(\sum_{n=1}^{N} w_n \log p_n\right),
        \]
        where $p_n$ is the modified $n$-gram precision, $w_n$ is the weight (uniform in our experiments), and BP is the brevity penalty.

        \item \textbf{chrF} measures character $n$-gram overlap using an F-score, which is more sensitive to morphological and subword differences. It is defined as:
        \[
        \text{chrF}_\beta = (1 + \beta^2) \cdot \frac{\text{Precision} \cdot \text{Recall}}{\beta^2 \cdot \text{Precision} + \text{Recall}},
        \]
        where precision and recall are computed over character $n$-grams, and we use $\beta = 2$ to emphasize recall.
\end{itemize}

\subsection{Variants of Transformer}
We study three architectural variants to understand the role of residual paths,
 normalization placement, and feed-forward networks.

\subsubsection{No Residual Connections}
We remove the residual connections in both the attention and feed-forward sublayers
while keeping all other settings unchanged.
This variant tests whether the shortcut paths are necessary for stable optimization and effective gradient flow in deep Transformer stacks.

Theoretically, without residual connections, the model may struggle to learn identity mappings,
 which can lead to vanishing gradients and hinder training convergence.
 We will empirically verify this by comparing the training dynamics and final performance of this variant against the baseline Transformer.

\subsubsection{Post-Norm Transformer}
We compare the common pre-norm design with a post-norm setting which is the original design in the "Attention is All You Need!" paper.
In the post-norm design, LayerNorm is applied after the residual connection, rather than before it as in the pre-norm design.
This means that the output of each sublayer is normalized after adding the input (residual) to the output of the sublayer,
rather than normalizing the sublayer output before adding the residual.
The post-norm structure can lead to different training dynamics and stability properties compared to pre-norm, and we will empirically compare their performance on our translation task.
The structure is illustrated in \Cref{fig:post_norm}.
\begin{figure}[H]
        \centering
        \includegraphics[width=0.8\linewidth]{figure/norm.png}
        \caption{Pre-norm and post-norm structures for a Transformer sublayer.}
        \label{fig:post_norm}
\end{figure}

\subsubsection{MoE with Linear Routing}
We introduce a lightweight Mixture-of-Experts (MoE) layer in the feed-forward block.
The router uses a simple linear projection to compute expert scores, and tokens are sent to top-$k$ experts without additional routing networks.

\begin{figure}[H]
        \centering
        \includegraphics[width=0.7\linewidth]{figure/moe.png}
        \caption{MoE feed-forward block with a linear routing projection.}
        \label{fig:moe}
\end{figure}

MoEs can increase model capacity without a proportional increase in computational cost, as only a subset of experts is activated for each token.
We will evaluate whether this simple MoE design can improve translation performance compared to the standard feed-forward network in the Transformer.

\section{Experiments}
\label{sec:experiments}

\subsection{Implement Details}
\label{sec:}


\subsection{Dataset}
\label{sec:dataset}
We will use the translation2019zh\cite{} Chinese to English translation dataset, which contains around 5200k sentence pairs for training and 40k sentence pairs for validation.
Here, we will use a tokenizer\cite{TiedemannThottingal:EAMT2020}to preprocess the dataset, and set the maximum sequence length to 192 for english sentences and 128 for chinese sentences.

\subsection{Experiment Result}

% baseline
% {"step": 550, "train_loss": 4.196592330932617, "val_loss": 4.252671337127685, "bleu": 0.1953611175372819, "chrf": 11.247890586049508, "eval_samples": 640.0}

% moe
% {"step": 600, "train_loss": 3.998141288757324, "val_loss": 4.094677042961121, "bleu": 0.17129140284159575, "chrf": 10.514405598699277, "eval_samples": 320.0}

% no residual
% {"step": 550, "train_loss": 5.750664710998535, "val_loss": 5.730744075775147, "bleu": 0.0, "chrf": 0.0, "eval_samples": 640.0}

% pre norm
% {"step": 550, "train_loss": 4.002425193786621, "val_loss": 4.052163410186767, "bleu": 0.14378664132134752, "chrf": 10.831988199882087, "eval_samples": 640.0}



\section{Conclusions}
\label{sec:conclusions}


\printbibliography

% \begin{appendices}

% \section{MATLAB Functions}
% Add your important MATLAB functions here with a brief implementation explanation. This is how to make an \textbf{unordered} list:
% \begin{itemize}
%     \item \texttt{y = linspace(x1,x2,n)} returns a row vector of \texttt{n} evenly spaced points between \texttt{x1} and \texttt{x2}.
%     \item \texttt{[X,Y] = meshgrid(x,y)} returns 2-D grid coordinates based on the coordinates contained in the vectors \texttt{x} and \texttt{y}. \text{X} is a matrix where each row is a copy of \texttt{x}, and \texttt{Y} is a matrix where each column is a copy of \texttt{y}. The grid represented by the coordinates \texttt{X} and \texttt{Y} has \texttt{length(y)} rows and \texttt{length(x)} columns.
% \end{itemize}

% \end{appendices}

\end{document}
